\chapter{Antecedentes}

Este capítulo ordena la literatura relevante en ocho ejes. Los tres
primeros provienen de la literatura argentina reciente: las limitaciones
de los enfoques contables (Rodríguez, 2023; Sanches, 2019), la economía
política del endeudamiento y de las reestructuraciones (Aruguete, 2021;
Roldán, 2021) y las consecuencias materiales del ciclo financiero
(Bohoslavsky y Cantamutto, 2024; Cantamutto et al., 2024). El cuarto
reúne la discusión internacional sobre la fatiga fiscal (Blanchard,
1990, 2019; Everaert y Jansen, 2018). Los cuatro restantes reúnen la
literatura internacional de simulación estocástica de deuda y tres
antecedentes específicos sobre Argentina: la tradición que combina
modelos VAR con simulación (Medeiros, 2012; Adler y Sosa, 2013;
Wijayanti y Rachmanira, 2020), las dos lecturas de la antesala de la
crisis de 2001-2002 (Cetrángolo et al., 1997; Calvo et al., 2003), la
controversia entre la ratio deuda/PIB y el balance patrimonial (Levy
Yeyati y Sturzenegger, 2023) y el costo real del \emph{default} (Hébert
y Schreger, 2017). El capítulo cierra con el área de vacancia.

\section{Primer eje: las limitaciones de los enfoques contables}

La literatura coincide en señalar las limitaciones de las metodologías
deterministas clásicas para evaluar la solvencia en economías
periféricas.

Rodríguez (2023) analiza la evolución de la deuda pública argentina
entre 2001 y 2022 desde el marco de la Nueva Economía Keynesiana.
Construye un indicador de ``efectos sobre la deuda pública'' que
descompone la dinámica de la ratio deuda/PIB en sus determinantes.
Encuentra que el indicador se ubicó en terreno negativo durante el 68\%
del período, lo que señala insostenibilidad y falta de margen fiscal, y
que la ratio deuda/PIB tomada de forma aislada no alcanza para concluir
sobre la sostenibilidad.

Sanches (2019), en su tesis de maestría, descompone la ratio deuda/PIB
argentina para aislar el peso del tipo de cambio, dado un perfil de
deuda con alta participación en moneda extranjera. Lo complementa con un
panel de países emergentes sobre la sensibilidad del diferencial de
rendimiento soberano a la composición por moneda de la deuda. Su
hallazgo central es que, con un perfil como el argentino, una
depreciación real afecta la ratio deuda/PIB tanto por el canal del
capital como por el de los intereses.

\section{Segundo eje: reestructuración y actores internacionales}

Aruguete (2021) examina la crisis de balanza de pagos de 2020, originada
en el sobreendeudamiento iniciado cuatro años antes, y la
reestructuración de la deuda en moneda extranjera habilitada por la Ley
27.541 de emergencia pública. Su diagnóstico subraya que los modelos de
sostenibilidad no pueden ignorar el uso macroeconómico de los fondos
prestados ni la vulnerabilidad que genera un perfil de vencimientos
concentrado en moneda extranjera.

Roldán (2021), en su tesina de grado, analiza la renegociación de 2020
con los tenedores privados de títulos bajo legislación extranjera desde
la teoría de juegos. Destaca el papel de las cláusulas de acción
colectiva, incorporadas a los nuevos títulos tras los litigios de 2014,
que obligan a los acreedores a aceptar el acuerdo una vez superado
cierto umbral de adhesión y reducen el poder de negociación de los
acreedores que no entran al canje (\emph{holdouts})\footnote{Acreedores
  que rechazan una reestructuración y reclaman judicialmente el pago
  completo.}.

Para esta tesis, estos trabajos importan porque ubican el origen de la
última reestructuración en el perfil de vencimientos en moneda
extranjera, el tramo de la deuda que la identidad de acumulación del
Capítulo 3 separa del tramo en pesos.

\section{Tercer eje: la tensión material y la restricción externa}

Bohoslavsky y Cantamutto (2024) sostienen que el sobreendeudamiento
argentino guarda una relación estructural con la especialización
productiva extractiva del país. Un criterio de sostenibilidad
restringido al balance fiscal corriente omite, por lo tanto, la
dimensión estructural de largo plazo. En la misma línea, Cantamutto et
al.~(2024) proponen leer la sostenibilidad desde la ``sostenibilidad de
la vida'': cuestionan que la capacidad de pago frente a los acreedores
pueda evaluarse con independencia de sus efectos sobre la reproducción
social. El planteo sistematiza un argumento previo de Cantamutto y Féliz
(2021). Para esta tesis, el eje respalda tratar la restricción externa
como contexto del análisis (sección 3.1) y recuerda que un superávit
primario sostenido tiene costos sociales que el modelo no mide.

\section{Cuarto eje: la discusión sobre la fatiga fiscal}

La fatiga fiscal es la hipótesis de que la respuesta del resultado
primario se debilita cuando la deuda es alta (Ghosh et al., 2013). Junto
a los trabajos que la documentan, la literatura incluye otros que
relativizan su urgencia. Blanchard (1990) ya propuso complementar la
ratio deuda/PIB con otros indicadores fiscales en lugar de tratarla como
un límite autosuficiente. Tres décadas después, Blanchard (2019)
argumenta que, cuando la tasa de interés real se ubica sistemáticamente
por debajo del crecimiento, el gobierno puede sostener déficits
primarios moderados sin que la ratio diverja.

Everaert y Jansen (2018) discuten la fatiga fiscal desde otro ángulo.
Retoman el panel de países de la OCDE (1970-2014) en el que se apoya
buena parte de esta literatura y testean si el patrón no lineal que
reportan Ghosh et al.~(2013) es una regularidad genuina o un efecto de
agregar países heterogéneos en un solo panel. Usan el estimador de
grupos medios (Pesaran y Smith, 1995), que permite una pendiente de
reacción fiscal distinta para cada país, y encuentran heterogeneidad
significativa: la fatiga fiscal es el promedio de comportamientos muy
distintos entre países y no un rasgo compartido por todos. La
implicancia para un análisis de un solo país, y en particular de
Argentina, es directa. Una función de reacción estimada sobre economías
avanzadas no es transferible sin más, porque esos parámetros provienen
de países con marcos de disciplina institucional que Argentina no tiene.
Esta advertencia respalda el diseño de caso único de esta tesis.

\section{Quinto eje: simulación estocástica de deuda y marcos de riesgo
soberano}

La literatura empírica sobre sostenibilidad de deuda soberana se
desplazó, en las últimas dos décadas, de las proyecciones deterministas
de un único escenario hacia simulaciones estocásticas que generan
distribuciones de probabilidad sobre la trayectoria de la deuda. El
origen de este giro en los organismos multilaterales son los gráficos de
abanico desarrollados en el FMI por Celasun et al.~(2006) y el trabajo
de Medeiros (2012) en la Comisión Europea. Medeiros combinó un modelo
VAR con una función de reacción fiscal de panel para quince países de la
Unión Europea y simuló miles de trayectorias de deuda/PIB. Su función de
reacción incluyó un término cúbico de la deuda rezagada, tomado de la
versión de documento de trabajo de Ghosh et al.~(2013), y encontró que
la respuesta del superávit primario se debilita a partir de una deuda de
80\% a 90\% del PIB. Esa noción de fatiga fiscal la absorbieron luego
los marcos de evaluación de solvencia europeos (Bouabdallah et al.,
2017). El FMI consolidó el enfoque probabilístico en su Marco de Riesgo
Soberano y Sostenibilidad de la Deuda (SRDSF; FMI, 2021), que combina el
gráfico de abanico con herramientas de realismo sobre el resultado
primario y que aplicó a Argentina en sus revisiones del programa (FMI,
2024, 2025).

Un segundo cuerpo de literatura llevó el problema a economías emergentes
expuestas a \emph{shocks} externos. Adler y Sosa (2013) estimaron, país
por país, modelos VAR para economías de América Latina con el
crecimiento mundial, el índice VIX y los precios de las materias primas
como variables externas, y simularon la trayectoria de la deuda bajo
escenarios externos adversos. Entre las variables globales, el VIX fue
la única con efectos significativos y robustos, de modo que la
sostenibilidad de la deuda en la región depende en buena medida de
factores de empuje externos (\emph{push})\footnote{Determinante externo
  al país receptor que impulsa los movimientos de capital, en contraste
  con los factores de atracción (\emph{pull}), asociados a condiciones
  domésticas.}. Adler et al.~(2016) estimaron un VAR de panel de los
flujos brutos de capital de 38 economías emergentes entre 1990 y 2012.
Los inversores locales cumplen un papel estabilizador ante los
\emph{shocks} de incertidumbre global y de la tasa larga de Estados
Unidos, aunque en el segundo caso la compensación es mucho más limitada.
Ante los \emph{shocks} de la tasa corta, residentes y no residentes
retiran capital de los emergentes por igual, con magnitudes
económicamente poco significativas.

La literatura de sistemas de alerta temprana (\emph{early warning
systems})\footnote{Modelos estadísticos que estiman la probabilidad de
  una crisis a partir de indicadores, sin pretender explicar el
  mecanismo causal.} buscó anticipar crisis de deuda con otro objetivo:
maximizar el poder predictivo, sin pretender explicar mecanismos.
Wijayanti y Rachmanira (2020) aplicaron el enfoque de señales de
Kaminsky et al.~(1998) y una regresión logística a un panel de 43 países
en desarrollo entre 1960 y 2017. Predijeron el 61,5\% de las crisis de
deuda con dos años de anticipación. El resultado expone una tensión:
estos sistemas maximizan el poder predictivo sobre un panel amplio,
mientras que el gráfico de abanico del FMI busca representar la
distribución de riesgo de un país en particular, y ambos enfoques pueden
dar señales distintas frente a un mismo episodio.

\section{Sexto eje: dos lecturas convergentes de la crisis de 2001-2002}

Dos trabajos centrados en Argentina, con instrumentos casi opuestos,
llegaron a diagnósticos convergentes en la antesala de la crisis de
2001-2002. Cetrángolo et al.~(1997) reconstruyeron el esquema de ahorro
e inversión del Sector Público Nacional no Financiero para 1975-1996.
Desagregaron el presupuesto en ocho rubros de ingresos y siete de gasto,
vincularon cada uno a su variable macroeconómica con coeficientes fijos
y aplicaron el filtro de Hodrick y Prescott (1997) para calcular un
déficit ajustado por el ciclo y un impulso fiscal, siguiendo a Blanchard
(1990). Su conclusión fiscal fue tranquilizadora: el riesgo fiscal
cíclico bajo la Convertibilidad era pequeño (0,3-0,6\% del PIB). Pero su
propia proyección para 1997-2001 mostraba un déficit de cuenta corriente
que superaba el 7\% del PIB hacia 2001, y advertía así que el riesgo
principal era externo, no fiscal.

Calvo et al.~(2003) llegaron a la misma advertencia con un instrumental
distinto: contabilidad de balanza de pagos, calibración de elasticidades
de demanda de bienes no transables y la identidad de dinámica de la
deuda. Mostraron que, ante una interrupción súbita del financiamiento
como la que siguió a la crisis rusa de 1998, el ajuste del tipo de
cambio real necesario para restablecer el equilibrio externo dependía de
la apertura comercial y del descalce de monedas de los pasivos. En ambos
rasgos, Argentina se ubicaba en el peor extremo entre los casos
comparados (Argentina, Brasil, Chile, Colombia y Ecuador). Uno de los
trabajos detecta el problema por el lado del flujo (cuenta corriente) y
el otro por el del \emph{stock} (descalce patrimonial). Ambos coinciden
en que el diagnóstico fiscal de corto plazo no bastaba para juzgar la
sostenibilidad.

\section{Séptimo eje: el balance patrimonial frente a la ratio
deuda/PIB}

Dos décadas después, la discusión volvió sobre si la ratio deuda/PIB es
una medida confiable de riesgo cuando hay un salto cambiario. Rodríguez
(2023) atribuye la caída de la ratio entre 2003 y 2011 a un efecto
contable de la apreciación real, más que a un desendeudamiento genuino.
El mismo mecanismo, en sentido inverso, explica que el salto cambiario
de fines de 2023 llevara la deuda al 156,7\% del PIB\footnote{El FMI
  revisó la cifra a 155,4\% en su informe de 2025 (FMI, 2025, anexo II,
  cuadro 2).}. En el marco del FMI, ese salto ensanchó el gráfico de
abanico hasta 122,7 puntos y llevó la señal mecánica de riesgo a
``alto'', que el propio personal técnico del organismo descartó por no
reflejar el riesgo de mediano plazo (FMI, 2024). Una vez disipado el
efecto de la devaluación, el ancho del abanico bajó a 74,2 puntos y la
señal volvió a ``moderado'' (FMI, 2025, anexo II).

Levy Yeyati y Sturzenegger (2023) llevaron la discusión un paso más
allá: propusieron reemplazar la ratio deuda/PIB por el patrimonio neto
del sector público: el valor presente de los ingresos menos el de los
gastos y pasivos explícitos. Lo simularon con un
\emph{bootstrap}\footnote{\emph{Bootstrap}: método de remuestreo que
  aproxima la distribución de un estimador repitiendo la estimación
  sobre muestras generadas a partir de los propios datos.} de un VAR de
tres variables (PIB, tipo de cambio real y tasa internacional). Para
Argentina obtuvieron un patrimonio neto medio de 2,2 veces el PIB, sin
valores negativos en la distribución simulada. Una devaluación real del
20\% apenas lo modifica (de 2,225 a 2,214 veces el PIB), porque el mayor
costo en dólares de la deuda se compensa con más recaudación de bienes
transables y con la licuación de gastos en pesos. El resultado choca con
el cálculo tradicional que los mismos autores citan de Calvo et
al.~(2003): con el enfoque de descalce de monedas, una devaluación del
50\% eleva la ratio deuda/PIB argentina de 36,5\% a 50,8\%. Los autores
advierten que su resultado benigno depende de que Argentina redujo la
dolarización de su deuda después de 2001. Con la estructura de deuda de
2001, el mismo \emph{shock} habría producido un deterioro severo.

\section{\texorpdfstring{Octavo eje: el costo real del
\emph{default}}{Octavo eje: el costo real del default}}

El último antecedente discute el costo \emph{ex post} del evento que
estos marcos buscan evitar. Hébert y Schreger (2017) usaron el litigio
de la República Argentina con NML Capital como experimento natural: los
fallos judiciales alteraron la probabilidad de \emph{default} percibida
por el mercado con independencia de las noticias económicas del día. Con
identificación por heterocedasticidad (Rigobon, 2003; Rigobon y Sack,
2004), estimaron que un aumento de 10 puntos porcentuales en la
probabilidad de \emph{default}, medida con los seguros de incumplimiento
(CDS), redujo en alrededor de 6\% el índice de acciones argentinas. El
paso de una probabilidad de 40\% a una de 100\%, lo efectivamente
ocurrido en 2014, implicó una pérdida de valor de mercado de las
empresas cotizantes de entre 30\% y 45\%, concentrada en exportadoras y
filiales de empresas extranjeras. Este trabajo cuantifica el costo del
evento extremo que los marcos estocásticos (Medeiros, 2012; FMI, 2021)
buscan mantener con baja probabilidad. Para esta tesis, ese costo
explica por qué interesa medir la probabilidad de que la deuda no se
estabilice, aunque la tesis no estima la probabilidad de un
\emph{default}.

\section{Área de vacancia}

La literatura revisada deja tres aspectos sin cubrir de manera conjunta
para el caso argentino.

\textbf{Identificación estructural de los canales de transmisión.} La
literatura local, desde el enfoque contable (Rodríguez, 2023; Sanches,
2019) y desde la economía política (Aruguete, 2021; Bohoslavsky y
Cantamutto, 2024; Cantamutto et al., 2024), vincula el esfuerzo fiscal y
el riesgo soberano de forma argumentativa, sin estimar un sistema en el
que los \emph{shocks} se propaguen entre el resultado fiscal, el riesgo
país, el tipo de cambio real y la actividad. La literatura internacional
modela esa transmisión, en general sobre paneles de países (Medeiros,
2012; Adler et al., 2016). Algunos trabajos incluyen a Argentina.
Celasun et al.~(2006) aplican su gráfico de abanico a una muestra que
contiene a Argentina, con un VAR en forma reducida, sin \emph{shocks}
identificados, y una función de reacción fiscal estimada en panel. Levy
Yeyati y Sturzenegger (2023) estiman un VAR para Argentina, pero la
variable que simulan es el patrimonio neto del sector público. Canova
(2005) identifica los \emph{shocks} de Estados Unidos y estudia su
transmisión a América Latina, sin llevarlos a la dinámica de la deuda.
En la literatura revisada no se encontró una estimación para Argentina
como caso único, sobre el perímetro consolidado del Tesoro y el BCRA,
cuyo resultado se exprese en la métrica de no estabilización del FMI.

En la evaluación más reciente de funciones de reacción fiscal para
América Latina y el Caribe, Argentina no integra el grupo de países con
series aptas para una estimación individual (González y Hernández,
2023). Esta tesis atiende esa objeción con frecuencia trimestral, que da
111 observaciones en la muestra de estimación (1998T2-2025T4), y con
variables binarias para los quiebres estructurales; la longitud y la
inestabilidad de las series argentinas siguen siendo, aun así, una
limitación del diseño (sección 7.6). El gráfico de abanico del FMI, por
su parte, aplica al escenario base \emph{shocks} tomados de la historia
de los determinantes de la deuda, sin identificar su origen (FMI, 2021,
párr. 52), y por eso no permite atribuir el riesgo a una perturbación
concreta.

\textbf{Un \emph{shock} externo identificado.} Los efectos de la
política monetaria de Estados Unidos sobre las economías emergentes se
estudiaron sobre los \emph{spreads}\footnote{\emph{Spread}: diferencia
  de rendimiento entre un bono y otro de referencia; en el caso
  soberano, entre los bonos del país y los del Tesoro de Estados Unidos.}
y la actividad (Uribe y Yue, 2006), sobre las variables macroeconómicas
de América Latina (Canova, 2005) y sobre los flujos de capital en
paneles de países (Adler et al., 2016). En la literatura revisada no se
encontró un trabajo que traduzca esos efectos en un cambio en la
probabilidad de que la deuda argentina no se estabilice. Levy Yeyati y
Sturzenegger (2023) incluyen la tasa internacional en su VAR, pero su
variable de interés es el patrimonio neto del sector público y su
sistema no incorpora el riesgo soberano ni la respuesta fiscal.

\textbf{El perímetro de la deuda en la estimación.} El marco del FMI
admite incluir los pasivos del banco central en la deuda pública. Las
letras de liquidez se incluyen cuando pueden financiar al gobierno,
cuando el gobierno responde de hecho por su servicio o cuando
representan un riesgo fiscal material, y el propio documento menciona
que en Argentina sumaban 5\% del PIB a fines de 2018 (FMI, 2021, anexo
II). Sin embargo, en su aplicación a Argentina, el FMI corre sus
herramientas sobre la deuda bruta del gobierno federal. La deuda
consolidada con el BCRA aparece solo como un cuadro complementario de
\emph{stocks}: 110,8\% del PIB en 2023 y 51,5\% en 2024, con pasivos del
BCRA que pasaron de 18,4\% a 2,2\% del PIB (FMI, 2025, anexo II, cuadro
2). En la literatura revisada no se encontró un trabajo que estime cómo
evolucionó esa deuda consolidada junto con sus determinantes a lo largo
del tiempo.

Esa omisión importa por razones institucionales. La autonomía formal que
el BCRA había obtenido en 1992 se redujo de manera progresiva desde
2002, y en especial entre 2010 y 2012, cuando la reforma de su Carta
Orgánica amplió los márgenes legales de asistencia al Tesoro (Damill et
al., 2015). En ese período, el Tesoro canceló deuda externa con reservas
del BCRA (Damill et al., 2015), y entre 2003 y 2015 las utilidades del
BCRA se convirtieron en una fuente regular de ingresos extraordinarios
del Tesoro (Porto, 2020). La absorción de esa liquidez se hizo con
pasivos remunerados: letras (LEBAC y NOBAC) hasta 2018, LELIQ y pases
entre 2018 y 2024. Es plausible que ese \emph{stock} haya respondido a
los mismos \emph{shocks} cambiarios y de riesgo que la deuda del Tesoro,
y esta tesis lo trata como una hipótesis. Un modelo estimado solo sobre
la deuda del Tesoro omite esa parte del endeudamiento y puede subestimar
la sensibilidad de la deuda a esos \emph{shocks}. Esto sigue siendo
cierto aunque, desde 2024, los pasivos remunerados del BCRA se hayan
reducido a niveles mínimos. Los pases pasivos se suspendieron en julio
de 2024 y su \emph{stock} pasó al Tesoro como Letras Fiscales de
Liquidez (LEFI), que el BCRA dejó de ofrecer en julio de 2025 (BCRA,
2025). A fin de 2025 el \emph{stock} de LEFI era nulo, frente a 11,1
billones de pesos un año antes (BCRA, 2026). La proyección parte de un
BCRA casi sin pasivos, pero los parámetros que la gobiernan se estiman
sobre tres décadas en las que no fue así.

Esta tesis atiende las tres omisiones con un solo diseño: un SVARX de
caso único para Argentina (1996-2025), estimado sobre la deuda
consolidada del Tesoro nacional y el BCRA, con identificación
estructural de los canales y la tasa de la FED como variable exógena.
Sus resultados alimentan una simulación de Monte Carlo que expresa el
riesgo con las métricas del propio FMI. La metodología se elige porque
reproduce la arquitectura que el FMI usa para evaluar la sostenibilidad,
la de Celasun et al.~(2006): un VAR que genera los \emph{shocks} y una
función de reacción fiscal que responde a la deuda. Aplicada con datos
argentinos y con la métrica de no estabilización del FMI, el resultado
de la tesis es comparable con la probabilidad que el organismo informa
para Argentina, y las diferencias pueden atribuirse al perímetro de la
deuda y a la identificación de los \emph{shocks}. Las variables se
eligen porque son los impulsores de la deuda en la identidad del FMI
(crecimiento, tipo de cambio real, costo del financiamiento a través del
riesgo país y resultado primario), a los que se suma la tasa de la FED
como el \emph{shock} externo que la literatura revisada vincula con esos
canales.
